\chapter*{ประวัติผู้เขียน}
\addcontentsline{toc}{chapter}{ประวัติผู้เขียน}

\begin{center}
    \vspace*{0.5cm}
    {\Large\bfseries ผู้ช่วยศาสตราจารย์ ดร.ชีวะ ทัศนา\par}
    \vspace{0.2cm}
    {\normalsize สาขาวิชาฟิสิกส์ คณะวิทยาศาสตร์และเทคโนโลยี มหาวิทยาลัยราชภัฏรำไพพรรณี\par}
    {\normalsize อีเมล: \texttt{chewa.t@rbru.ac.th}\par}
    \vspace{0.8cm}
\end{center}

\noindent\textbf{ประวัติการศึกษา}
\begin{itemize}[leftmargin=2.0em, itemsep=4pt]
    \item ปร.ด. (ฟิสิกส์ประยุกต์), มหาวิทยาลัยเทคโนโลยีสุรนารี
    \item วท.ม. (ฟิสิกส์ประยุกต์), มหาวิทยาลัยเชียงใหม่
    \item วท.บ. (ฟิสิกส์ คอมพิวเตอร์และอิเล็กทรอนิกส์), มหาวิทยาลัยราชภัฏรำไพพรรณี
\end{itemize}

\vspace{0.6cm}
\noindent\textbf{ความเชี่ยวชาญและงานวิจัยที่สนใจ}
\begin{itemize}[leftmargin=2.0em, itemsep=4pt]
    \item ดาราศาสตร์ฟิสิกส์และการสร้างแบบจำลองระบบดาวคู่ (Astrophysics \& Contact Binary Systems)
    \item ฟิสิกส์เชิงคำนวณและการจำลองทางคอมพิวเตอร์ (Computational Physics \& First-Principles Calculations)
    \item เทคโนโลยีความเป็นจริงขยายและการประมวลผลเชิงพื้นที่ (Spatial Computing, AR/VR/XR \& OpenXR)
    \item การออกแบบสื่อการเรียนรู้สะเต็มศึกษาและนวัตกรรมการศึกษาเชิงรุก (STEM Education Innovations)
\end{itemize}

\vspace{0.6cm}
\noindent\textbf{ผลงานทางวิชาการและงานวิจัยดีเด่น}
\begin{itemize}[leftmargin=2.0em, itemsep=4pt]
    \item ผลงานวิจัยตีพิมพ์ในวารสารระดับนานาชาติตามฐานข้อมูล Scopus / Web of Science ด้านการคำนวณโครงสร้างทางอิเล็กทรอนิกส์ด้วยทฤษฎีฟังก์ชันความหนาแน่น (Density Functional Theory) และฟิสิกส์ดาราศาสตร์
    \item ผู้พัฒนาโครงการแพลตฟอร์มการเรียนรู้เสมือนจริงและห้องปฏิบัติการวิทยาศาสตร์จำลองสามมิติ มหาวิทยาลัยราชภัฏรำไพพรรณี
    \item ผู้แต่งตำราและเอกสารคำสอนวิชาคณิตศาสตร์สำหรับฟิสิกส์ ฟิสิกส์แผนใหม่เชิงคำนวณ และการประมวลผลเชิงพื้นที่เพื่อการศึกษา
\end{itemize}
\clearpage
